\documentclass[12pt,a4paper]{article}
\usepackage[spanish]{babel}
\usepackage[utf8]{inputenc}
\usepackage[T1]{fontenc}
\usepackage{geometry}
\usepackage{booktabs}
\usepackage{longtable}
\usepackage{hyperref}
\geometry{margin=2.5cm}

\title{Informe sobre Virtualización de Infraestructura para Sistemas Distribuidos}
\author{Nombre(s) del/los integrante(s) \\\ Asignatura \\\ Docente}
\date{Fecha}

\begin{document}
\maketitle
\thispagestyle{empty}

\vspace{2cm}
\begin{center}
\textit{\large Plantilla de informe sin contenido. Reemplace los apartados con su trabajo.}
\end{center}

\newpage
\tableofcontents
\newpage

\section{Introducción}
\label{sec:introduccion}
% Breve presentación del tema y su importancia para los sistemas distribuidos.

\section{Conceptos fundamentales}
\label{sec:conceptos}
% Explicación de virtualización de servidores e infraestructura.

\subsection{Hipervisor}
% Definición del hipervisor.

\subsection{Máquina virtual}
% Definición de máquina virtual.

\subsection{Servidor físico o host}
% Definición del servidor físico o host.

\subsection{Clúster de virtualización}
% Definición del clúster de virtualización.

\subsection{Red virtual}
% Definición de red virtual.

\subsection{Almacenamiento virtualizado o compartido}
% Definición de almacenamiento virtualizado o compartido.

\subsection{Alta disponibilidad}
% Definición de alta disponibilidad.

\subsection{Infraestructura como servicio (IaaS)}
% Definición de IaaS.

\section{Relación con sistemas distribuidos}
\label{sec:relacion}
% Explicación de cómo la virtualización ayuda al funcionamiento de sistemas distribuidos.

\section{Herramientas y soluciones actuales}
\label{sec:herramientas}
% Descripción de las cuatro herramientas seleccionadas.

\subsection{Herramienta 1}
% Qué es, modelo de uso, propósito, ventaja y limitación.

\subsection{Herramienta 2}
% Qué es, modelo de uso, propósito, ventaja y limitación.

\subsection{Herramienta 3}
% Qué es, modelo de uso, propósito, ventaja y limitación.

\subsection{Herramienta 4}
% Qué es, modelo de uso, propósito, ventaja y limitación.

\section{Matriz comparativa simplificada}
\label{sec:matriz}
\begin{table}[ht]
\centering
\begin{tabular}{@{}llllll@{}}
\toprule
Herramienta o solución & Tipo de tecnología & Modelo de uso & Ventaja principal & Limitación o desafío & Uso recomendado en sistemas distribuidos \\
\midrule
% Rellene las filas con las herramientas seleccionadas.
\textit{Herramienta 1} & & & & & \\
\textit{Herramienta 2} & & & & & \\
\textit{Herramienta 3} & & & & & \\
\textit{Herramienta 4} & & & & & \\
\bottomrule
\end{tabular}
\caption{Matriz comparativa de herramientas y soluciones de virtualización.}
\label{tab:matriz-comparativa}
\end{table}

\section{Usos en contextos reales}
\label{sec:usos}
% Análisis de dos contextos reales.

\subsection{Contexto real 1}
% Explicación del primer contexto real.

\subsection{Contexto real 2}
% Explicación del segundo contexto real.

\section{Tabla de usos en contextos reales}
\label{sec:tabla-usos}
\begin{table}[ht]
\centering
\begin{tabular}{@{}llllll@{}}
\toprule
Contexto real & Necesidad principal & Solución aplicable & Componentes distribuidos involucrados & Beneficio principal & Riesgo o desafío \\
\midrule
% Rellene las filas con los dos contextos analizados.
\textit{Contexto 1} & & & & & \\
\textit{Contexto 2} & & & & & \\
\bottomrule
\end{tabular}
\caption{Tabla de usos de la virtualización en contextos reales.}
\label{tab:usos-contexto}
\end{table}

\section{Conclusiones}
\label{sec:conclusiones}
% Síntesis de lo aprendido y utilidad de la virtualización para infraestructura distribuida.

\section{Declaración de uso de herramientas de IA}
\label{sec:ia}
% Indicar si se usaron herramientas de IA, cuáles y para qué propósito.

\section{Referencias}
\label{sec:referencias}
\begin{itemize}
\item Fuente 1.
\item Fuente 2.
\item Fuente 3.
\item Fuente 4.
\end{itemize}

\end{document}
